\section{Algorithm for Conformal Subgraphs}
\label{sec:algorithm}

The previous section reformulated conformal compression as a combinatorial optimization problem on a weighted hypergraph~$H=(V,\mathcal E)$. Given an optimal (but unknown) vertex set~$O$ that covers all but an $\epsilon$~fraction of the total edge weight our goal is to design an efficient algorithm that finds a small vertex subset whose induced subgraph captures nearly the same total weight.

While the general Densest $k$-Subgraph problem is hard to approximate, we show that the conformal compression regime---where the target subgraph must contain the vast majority of the total weight---allows for efficient constant-factor approximations. We further show that the resulting algorithm satisfies monotonicity (Definition~\ref{def:monotone}).  %We present an algorithm based on a linear programming (LP) relaxation. 

\subsection{LP Rounding Algorithm} 
\label{sec:lp}
We first formulate the continuous relaxation of the problem. Let $x_v \in [0,1]$ indicate the inclusion of vertex $v$, and $z_e \in [0,1]$ indicate the inclusion of hyperedge $e$. Since a hyperedge is induced only if all its vertices are chosen, we enforce $z_e \le x_v$ for all $v \in e$.

\begin{equation}
\label{eq:lp}
\begin{aligned}
\text{minimize} \quad & \sum_{v \in V} x_v \\
\text{subject to} \quad & z_e \le x_v, \quad \forall e \in \mathcal{E}, \forall v \in e \\
& \sum_{e \in \mathcal{E}} w_e z_e \ge (1-\epsilon)W \\
& 0 \le x_v, z_e \le 1.
\end{aligned}
\end{equation}

The algorithm proceeds as follows:
\begin{enumerate}
  \item Construct and solve the linear program (\ref{eq:lp}) to obtain an optimal fractional solution $(x^*, z^*)$.
  \item Fix a parameter $\kappa \in \left(0,\frac{1}{\epsilon}-1 \right)$ determining the trade-off between size and coverage.
%  \item Set the rounding threshold $\rho := \frac{\kappa}{1+\kappa}$.
  \item Return the set $K := \left\{ v \in V : x^*_v \ge \rho := \frac{\kappa}{1+\kappa} \right\}$.
\end{enumerate}


\paragraph{Approximation analysis.}
We now analyze the approximation guarantees of the LP rounding scheme. Let $r = |O|$ be the size of the optimal solution. %Note that the integer assignment corresponding to $O$ (setting $x_v=1$ for $v \in O$ and $z_e=1$ for induced edges) is feasible for the LP. Therefore, 
Note that the objective value of the LP satisfies $\sum_{v \in V} x^*_v \le r$.

\begin{theorem}[Proved in Appendix~\ref{app:mainproof}]
\label{thm:main-det}
For any $\kappa>0$, the algorithm above returns a set $K$ that is a $\Bigl(1+\kappa,\;1+\frac{1}{\kappa}\Bigr)$ bicriteria approximation: that is,
\[
W - e(K) \le (1+\kappa)\,\epsilon W
\qquad\text{and}\qquad
|K| \le \Bigl(1+\frac{1}{\kappa}\Bigr)\,r.
\]
%In the specific case of $\kappa=1$ (corresponding to threshold $\rho=1/2$), this yields a $(2,2)$-approximation.
\end{theorem}

In the distribution-free setting (Section~\ref{sec:free}), this yields an approximately optimal compression of $S_{d^*}(A^*)$ while relaxing $(1-\tau)$ by a constant factor.  We interpret the guarantee for the fixed context setting from Section~\ref{sec:fixed} in Section~\ref{sec:fixed_optimal} below. Finally, for the distribution-based setting in Section~\ref{sec:distr2}, Theorem~\ref{thm:main-det} directly yields an approximate optimality--coverage trade-off when $w_e$ capture the posterior probabilities of the ground truth, and this procedure can be conformalized via the learn-then-test method.

Note that though the inequality $W_{\text{lost}} := W - e(K) \le (1+\kappa)\epsilon W$ holds unconditionally for any graph, this bound is only meaningful in the conformal compression regime where the target subgraph retains a large fraction of edges (i.e., $\epsilon$ is a small constant). In the classical Densest $k$-Subgraph (DkS) regime, the target subgraph contains a sublinear fraction of the edges, meaning $\epsilon \to 1$. If we plug $\epsilon \approx 1$ into our guarantee, the bound becomes $W_{\text{lost}} \le (1+\kappa)W$, which is vacuous (it allows the algorithm to lose all edges). Therefore, while Theorem~\ref{thm:main-det} requires no assumptions, the utility of the bicriteria approximation is mainly in the high-coverage (linear) regime, which (quite surprisingly) bypasses classical DkS hardness barriers. 

\subsection{Parametric Min-cuts and Monotonicity} 
\label{sec:mon}
We now show that our algorithm yields valid conformal guarantees in the distribution-free setting. Specifically, we prove that the solutions generated by the LP rounding form a nested family of subgraphs as the target coverage $\tau$ increases, hence satisfying Definition~\ref{def:monotone}. 

We establish this result by analyzing the Lagrangian relaxation of the LP. This Lagrangian is identical to the Lagrangian for the densest ratio sub-hypergraph, which is shown in~\cite{ChekuriQT} to  map to a Parametric Minimum Cut problem~\cite{gallo1989fast}. The monotonicity of the Lagrangian now follows from~\cite{gallo1989fast}. Our subsequent contribution is to show that this monotonicity is preserved when we move from the Lagrangian to the LP optimum, and from the LP optimum to the rounding. Indeed, we show that solving the parametric minimum cut problem suffices to extract all the compressed subgraphs for various settings of the compression parameter, and these solutions are identical to those found by the LP rounding algorithm!

\paragraph{Lagrangian.} In more detail, consider the Lagrangian of the LP coverage constraint with a multiplier $\lambda \ge 0$. The objective becomes:
\begin{align*}
    \label{eq:lagrangian_lp}
    \mathcal{L}(\lambda) = \min_{x, z \in [0,1]} \left( \sum_{v \in V} x_v - \lambda \sum_{e \in \mathcal{E}} w_e z_e \right) \\ 
    \text{s.t.} \quad z_e \le x_v \quad \forall e \in \mathcal{E}, \forall v \in e.
\end{align*}

The above formulation is similar to the LP formulation for densest ratio subgraph in~\cite{charikar2000greedy}, and their proof implies the above formulation has an integer optimum (see Lemma~\ref{lem:integrality} for a proof). We can therefore interpret the optimization variables as defining a vertex set $K = \{v \in V \mid x_v = 1\}$. Due to the constraints $z_e \le x_v$, a hyperedge is induced ($z_e=1$) if and only if all its vertices are in $K$. Thus, for a fixed $\lambda$, the Lagrangian minimization problem is equivalent to finding a subgraph $K \subseteq V$ that minimizes:
\begin{equation}
    \Phi(K, \lambda) = |K| - \lambda \cdot W_{\text{induced}}(K).
\end{equation}

\paragraph{Parametric Min-cuts.} It is shown in~\cite{ChekuriQT} in the context of densest ratio sub-hypergraphs that $\mathcal{L}(\lambda)$ is precisely a standard Minimum $s-t$ Cut problem on a flow network. For this, we construct a flow network $D_\lambda$ with nodes $\{s, t\} \cup \{u_e \mid e \in \mathcal{E}\} \cup \{n_v \mid v \in V\}$ and the following edges: 

\begin{itemize}
    \item For each hyperedge $e$, an edge $(s, u_e)$ with capacity $C_{s, u_e} = \lambda w_e$.
    \item For each vertex $v$, an edge $(n_v, t)$ with capacity $C_{n_v, t} = 1$. 
    \item For each inclusion $v \in e$, an edge $(u_e, n_v)$ with capacity $C_{u_e, n_v} = \infty$.
\end{itemize}  

Let $K$ denote the vertices $n_v$ on the source side of the cut. Then the cut capacity is precisely
$ |K| + \lambda \cdot \left(\sum_{e \in \mathcal E}  w_e - W_{\text{induced}}(K)\right),$
which shows $\Phi(K,\lambda)$ is minimized by the minimum cut. 

In this graph, the set of nodes on the source side is exactly the union of the selected vertices and the induced hyperedges: $X(\lambda) = \{n_v \mid v \in K\} \cup \{u_e \mid e \text{ induced by } K\}$. The parameter $\lambda$ appears only in the capacities of the edges incident to the source $s$, and these capacities $C_{s, u_e}(\lambda) = \lambda w_e$ are non-decreasing functions of $\lambda$. This is a standard instance of the \textit{Parametric Minimum Cut} problem, whose output is the subgraph $X(\lambda)$.

%We now invoke  \cite{gallo1989fast} states that for such networks, the source-side set of the minimum cut forms a monotonic collection of sets with respect to the parameter $\lambda$.
A result of~\cite{gallo1989fast} for parametric min-cuts now directly implies the following lemma:

\begin{lemma}[Monotonicity of Lagrangian] 
\label{lem:lag-mon}
Let $\gamma$ denote the maximum size of any hyper-edge, and let $m = |\mathcal E|$ and $n = |V|$. Then in $\tilde{O}(\gamma (m+n)^2)$ time, we can compute a sequence of vertex subsets $\emptyset = S_0 \subset S_1 \subset S_2 \subset \cdots \subset S_k \subseteq V$ where $k \le n$, such that for any $\lambda \ge 0$, the solution $X(\lambda)$ is one of these subsets. Further, if $\lambda_1 < \lambda_2$, the optimal subgraphs satisfy $X(\lambda_1) \subseteq X(\lambda_2)$.
\end{lemma}

\paragraph{Algorithm.} Fix a slack parameter $\kappa$ as in the LP rounding algorithm.  Given the coverage parameter $\tau := 1- \epsilon$, the final algorithm is simple: Among the subgraphs $S_0, S_1, \ldots, S_k$ constructed above, set 
\[ \hat{K}_{\tau} = \operatorname*{argmin}_{S \in \{S_1, \dots, S_k\}} \left\{ |S| : W - W_{\text{induced}}(S) \le (1+\kappa)\epsilon W \right\}. \]

Note that by Lemma~\ref{lem:lag-mon}, we have $\hat{K}_{\tau_1} \subseteq \hat{K}_{\tau_2}$ for all $\tau_1 \le \tau_2$, so that this construction is monotone. Further, the set of all subgraphs can be constructed in $\tilde{O}(\gamma (m+n)^2)$ time, where $\gamma$ is the largest size of a hyperedge.

Our main result is the following surprising statement, showing this algorithm is identical to the LP rounding algorithm:

\begin{theorem}[Proved in Appendix~\ref{app:monotonicity}]
\label{thm:mon}
Fix a slack parameter $\kappa$. For any target coverage $\tau  = 1 -\epsilon$, let $K_\tau$ be the discrete subgraph obtained by the LP rounding algorithm, and $\hat{K}_{\tau}$ be the subgraph obtained by the parametric min-cut algorithm. Then,  $K_{\tau} = \hat{K}_{\tau}$, so that both algorithms yield the bicriteria approximation guarantee in Theorem~\ref{thm:main-det} and produce subgraphs that are monotone in $\tau$.
\end{theorem}


%Indeed, Theorem~\ref{thm:mon} also implies the following corollary, showing the {\em entire} sequence of subgraphs for different thresholds $\tau$ can be efficiently computed.

%\begin{corollary}[Proved in Appendix~\ref{app:monotonicity}]
%\label{cor:nested}
%Let $\gamma$ denote the maximum size of any hyper-edge, and let $m = |\mathcal E|$ and $n = |V|$. Then in $\tilde{O}(\gamma (m+n)^2)$ time, we can compute a sequence of vertex subsets $\emptyset = S_0 \subset S_1 \subset S_2 \subset \cdots \subset S_k \subseteq V$ where $k \le n$, with the following property: For any target coverage $\tau = 1-\epsilon$ and slack $\kappa$, the subgraph $K_{\tau}$ output by LP rounding  corresponds to the smallest set in this sequence satisfying the bicriteria constraint:
%\[ K_{\tau} = \operatorname*{argmin}_{S \in \{S_1, \dots, S_k\}} \left\{ |S| : W - W_{\text{induced}}(S) \le (1+\kappa)\epsilon W \right\}. \]
%\end{corollary}

\paragraph{Number of Hyperedges.}
\label{sec:poly-m-remark}
In our algorithms, we assumed the number of hyperedges \(|\mathcal{E}|\) is polynomially bounded in the number of vertices \(|V|\).  This is a natural assumption in many applications and it simplifies both presentation and implementation of the algorithm.  In Appendix~\ref{sec:sampling}, we show that even when the full hyperedge set output by the Pre-Filtering Stage 1 (the set $S_{d^*}(A^*)$) is exponentially large, we can replace this set by a polynomial-size random sample via the uniform convergence theorem~\cite{vapnik}, approximately preserving our guarantees.  
%This step requires only an oracle that draws hyperedges i.i.d.\ proportional to their weights/scores \(w_e\), and 
For applications like navigation, classification, and trip planning, we develop efficient sampling oracles for $S_{d^*}(A^*)$ under canonical distance functions $f(A,B)$. Uniform convergence is also relevant to our approximate optimality guarantees in the fixed context framework  below.

\subsection{Optimal Compression in Fixed Context}
\label{sec:fixed_optimal}
We now show how Theorem~\ref{thm:main-det} can also be used to provide (approximate) size-optimality guarantees in the distribution-free conformal prediction framework in the fixed-context setting described in Section~\ref{sec:fixed}; note that this also captures the setting of our experiments on Navigation in Section~\ref{sec:nav_experiments}. In this setting, the hypergraph $G(V,\mathcal E)$ is fixed and we observe $i.i.d.$ hyperedge samples $Y_1,Y_2,\ldots,Y_T$ from some unknown distribution. For given $\phi \in (0,1)$, the goal is to output $K^*$, the subgraph satisfying the conformal guarantee $\mathbb{P}(Y_{T+1} \in K^*) \ge \phi$, while minimizing $|K^*|$. 

We adapt the split conformal framework in~\cite{gao2025volume} to derive an oracle inequality for compression. 
Our algorithm will ensure the coverage guarantee (Validity) under just exchangeability of the samples $Y_1, \dots, Y_{T+1}$. If in addition, the samples are i.i.d. and the number of samples is large enough, then we also get an (approximate) size optimality guarantee (Compression) by combining Theorem~\ref{thm:main-det} with uniform convergence (Lemma~\ref{lem:uniform-sample}).  Our algorithm requires a monotone compression algorithm that for any coverage parameter $\phi$, finds a subgraph whose coverage is close to that parameter. We adapt the procedure in Section~\ref{sec:mon} for this as follows: 

\begin{enumerate}
    \item Use half of the samples, $\{Y_j\}_{j=1}^{T/2}$, to find the nested sequence $S_0 \subset S_1  \subset \cdots \subset S_k$  as in Lemma~\ref{lem:lag-mon}.
    \item Find the subgraph $S_i$ whose coverage on the other half of the samples, $\{Y_j\}_{j=T/2+1}^{T}$, is at least $\phi$. %$1 - \hat{\phi} = (1+\kappa)\cdot(1-\phi)$. 
    If $S_k$ has larger mis-coverage, add vertices to it in some fixed order till the coverage condition is satisfied.
    \item Iteratively delete the vertices in $S_i \setminus S_{i-1}$ in some fixed order till the coverage becomes close to $\phi$. 
\end{enumerate} 

%This yields the following set of guarantees: %from any such monotone algorithm, and it provides a compression-optimality guarantee provided the number of samples is sufficiently large relative to the VC-dimension, which in our case is at most $|V|$. (See the uniform convergence result in Appendix~\ref{sec:sampling}.)  This yields the following:

\begin{theorem}[Proved in Appendix~\ref{app:fixed}]
\label{cor:efficiency}
In the fixed context setting, the above procedure yields a subgraph $\hat{K}$ satisfying:
\begin{enumerate}
    \item \textbf{Validity:} If $Y_1, Y_2, \ldots, Y_{T+1}$ are exchangeable, then 
    $$\mathbb{P}(Y_{T+1} \in \hat{K}) \ge \phi.$$
    \item \textbf{Approximate Compression (restates Theorem~\ref{thm:main-det}):} Let $Y_1, Y_2, \ldots, Y_{T+1} \sim \mathcal{D}$ be $i.i.d.$ random variables. For constants $\gamma, \kappa > 0$, let 
    $$K^* = \mbox{argmin}_{K} |K| \ \  \mbox{s.t.} \ \  \mathbb{P}(Y_{T+1} \in K) \ge \hat{\phi} + \gamma,$$  
     where $1 - \hat{\phi} = \frac{1-\phi}{1+\kappa}$. If $T = \Omega\left(\frac{|V| }{\gamma^2}\right)$, then 
    $$|\hat{K}| \le \left(1+\frac{1}{\kappa} + o(1) \right) |K^*|.$$
\end{enumerate}
\end{theorem}

Note that the algorithm does not depend on $\kappa$, so that the compression guarantee holds simultaneously for all $\kappa > 0$. Further note that for a given $\phi$, the procedure in~\cite{gao2025volume} uses an arbitrary greedy method to construct a monotone collection of subgraphs. Our construction via Theorem~\ref{thm:mon} is more robust when the sample size $T$ is much smaller than the uniform convergence threshold, so that a theoretical bound on compression is not possible. This aspect is particularly relevant to our experiments in Appendix~\ref{sec:real} on real-world data that can potentially be sparse.

%\paragraph{Feasibility.} We briefly note that since we assume $\epsilon < 1/2$, the set $K$ is guaranteed to be non-empty for reasonable $\kappa$. Specifically, for the balanced case $\kappa=1$ ($\rho=1/2$), if $K$ were empty, then all $x^*_v < 1/2$, implying all $z^*_e < 1/2$. This would yield a total fractional weight $< W/2$, contradicting the constraint $(1-\epsilon)W > W/2$. Thus, the algorithm always returns a non-trivial conformal subgraph.


